\FloatBarrier
\section{Hub cover sharing with holes retained}
\label{sec:hub-covers}

Upcoming in GEMC2 6.2 and CLAS12 GEMC Systems 0.4, Plate1 through Plate6 are renamed hubCover1 through
hubCover6. The plain hubCover1 mesh supplies copies 2, 4, and 5; the holed hubCover3 mesh supplies copy 6.
Both canonical STLs are byte-identical to their engineering originals. All holes, thicknesses, contours,
G4\_Al material, colors, and six physical placements are retained in both owning repositories and their
runtime collections.

The four plain exports each contain eight vertices and 12 triangles in one closed manifold component of genus
zero. The two holed exports each have 75 vertices and 150 triangles, with genus one. All have zero boundary
edges and no detected self-intersections. Nominal sector rotations match every vertex one to one: the largest
plain-cover residual is 0.000005540 mm, and the holed pair agrees to numerical roundoff. Three plain exports
use the alternate diagonal on one planar quadrilateral. The builder verifies identical perimeters, four
corners, and coplanarity, establishing surface equivalence despite different triangle connectivity.

The two sources already have simple tessellations. Optimization therefore shares solids without further
decimation or hole filling. Every canonical triangle remains, including the complete hole wall and rim in
hubCover3. Passive copies use exact integer sector angles and positions with three decimal places in
millimetres. No contact relief or empirical translation is introduced.

Unique cover facets fall from 348 to 162, saving 186 (53.45\%). Active cover STL storage falls from 17,904 to
8,268 bytes. The six placements still contain 348 placed facets. Total signed aluminum volume is
7,346,686.175700 cubic millimetres before and after, with a relative change below one part in a trillion from
numerical roundoff. The complete magnets collection now contains 24 STL sources, 69,564 unique facets, and 39
CAD placements. GEMC3 keeps sources and copies in magnets\_cad; GEMC2 uses 24 CAD sources and 15 native
magnets\_copies rows per inherited run, including the existing jacket and coil copies.

All nine GEMC3, GEMC2 SQLite/CAD, and GEMC2 TEXT/CAD construction checks pass across runs 11, 3029, and 4763.
Semantic comparison matches all three native solenoid configurations and the three 39-row CAD banks. Unrelated
native geometry, materials, and DC banks are unchanged at synchronization. The builder reproduces both
canonical STLs, the health report, and all three catalogs byte for byte. Original engineering exports remain
archived.

The dense run-11 fixture contains native DC, both DC CAD groups, solenoid, magnets, and native/CAD beamline:
139 volumes including the world. Readout is disabled. Renaming volumes changes Geant4 construction/check
ordering and thus assigns surface samples differently among unchanged volumes. The initial comparison records
four added and four removed ordered contact pairs, all among unchanged non-cover volumes. A controlled
baseline keeps the six independent original meshes and original poses but uses the final hubCover names. With
seed 12345 and 100,000 surface samples per tested volume, this baseline and the final shared geometry report
exactly the same 22 ordered overlap pairs. No new sampled overlap pair is introduced. Existing contacts,
including covers against the upstream jacket, and orientation warnings remain. Sampling does not establish
mathematical clearance.

The timing fixture isolates the six cover placements at run 11. The baseline builds six independent solids,
and the final fixture builds two solids with four copies. A fixed executable reports geant4-11-03-patch-02
[MT]   (25-April-2025). It uses the serial run manager, FTFP\_BERT, default production cuts, and one execution
thread. Field, optical physics, GUI, readout, and event output are disabled. Five-GeV/c electrons, protons,
positive pions, and neutrons start at the origin, uniform in solid angle from 1.5 to 5.5 degrees with full
azimuth, covering these small-radius parts. There is no target, luminosity normalization, or pileup. Five
repetitions alternate fixture order and use 2,000 identical LUND primaries per particle at seeds 12345--12349.
Ten zero-event processes, forty tracking processes, and two isolated overlap checks are scheduled. Idle-sleep
prevention remains enabled throughout. The executable hash is checked at both ends. Inputs are frozen
independently of concurrent framework builds and experiment database regeneration.

One of the 52 scheduled processes fails. Raw logs, exception codes, and inputs are retained. The original
proton
fixture at seed 12345 raises PART122 because Na24 is already registered in the Geant4 particle table. Tracking
statistics include only pairs where both versions complete all 2,000 events; failures exclude both members of
that seed pair. These conditional measurements do not establish an unconditional tracking-performance gain.

Median zero-event process wall time changes from 0.13304 to 0.13308 s, a before/after ratio of 1.000. Median
CPU time changes from 0.090 to 0.090 s, and peak RSS from 58.05 to 57.69 MiB. Paired startup ratios range from
0.997 to 1.392. Tracking entries below are median internal transport wall times. Ratios above one indicate
faster execution; paired ranges show seed dependence, and ranges spanning one do not establish a consistent
gain in this sample. The principal demonstrated benefit is fewer unique solids and facets. These isolated
measurements do not establish a full-CLAS12 speedup or detector-response equivalence.

\begin{table}[htbp]
\centering
\setlength{\tabcolsep}{6pt}
\caption{Hub-cover tracking; completed pairs of 2,000-event processes only.}
\label{tab:hub-covers}
\begin{tabular}{lrrrrr}
\toprule
Particle & Pairs & Before (s) & After (s) & Ratio & Paired ratio range \\
\midrule
Electron & 5 & 4.45213 & 4.47410 & 0.995 & 0.967--1.108 \\
Proton & 4 & 4.46378 & 4.67747 & 0.954 & 0.904--1.018 \\
Positive pion & 5 & 4.01745 & 3.93746 & 1.020 & 0.881--1.119 \\
Neutron & 5 & 2.44995 & 2.73761 & 0.895 & 0.759--1.203 \\
\bottomrule
\end{tabular}
\end{table}

\begin{figure}[htbp]
\centering
\includegraphics[width=\linewidth]{figures/hub_cover_comparison.pdf}
\caption{
Original and shared hub-cover footprints coincide in all six sectors (left). The hole rim retained by
hubCover3 and its hubCover6 copy is unchanged (right); all 150 canonical facets remain. Plain covers 1, 2, 4,
and 5 retain their original 12-triangle shape.}
\label{fig:hub-covers}
\end{figure}

data/magnets\_hub\_covers and magnets\_hub\_covers\_archive.zip retain mesh equivalence and topology checks,
facet and volume metrics, all nine construction logs, semantic comparisons, the controlled and initial dense
overlap fixtures, repeated timings, matched primaries, raw logs, and the fixed executable hash. The archive
includes original exports, final collected geometry and meshes, catalogs, builders, benchmark, validation
code, and figure-generation inputs. The executable is not redistributed; previous study archives remain
unchanged.
